Class-imbalance corrections change the training prevalence of clinical risk models and may distort predicted risks. We run a small, reproducible CPU study on the breast cancer Wisconsin data bundled with scikit-learn, at natural prevalence and with training positives subsampled to 1:5, 1:20 and 1:50. We compare no correction, class reweighting, random oversampling (ROS), SMOTE (implemented directly) and decision-threshold moving, for logistic regression (LR) and gradient boosting (GB), over 10 random splits with prevalence-matched weighted evaluation. The result is mostly null on discrimination. For LR, AUROC differs from no correction by at most 0.003 in any cell, and at 1:5 reweighting and ROS are slightly worse in AUROC, AUPRC and Brier score (unadjusted paired $p\approx0.03$--$0.05$). For GB the differences are variable; only ROS at 1:5 shows an AUPRC gain above 0.01 with unadjusted $p<0.05$. Brier score is higher with a correction in 11 of 12 cells at 1:20 and 1:50, but never significantly. The clearest calibration effect is a shift of risk: for LR every correction lowers the calibration intercept in all 12 cells ($p\le0.0002$), turning underestimation into overestimation at 1:5, 1:20 and 1:50, while the absolute intercept is not consistently larger. At the 0.5 cut-off all three corrections raise LR sensitivity at every prevalence and LR balanced accuracy at 1:5, 1:20 and 1:50; threshold moving, which leaves every probability unchanged by construction, raises both more at 1:20 and 1:50. For GB neither helps reliably at 1:20 and 1:50. The analytic prior-shift logit offset does not consistently bring the Brier score closer to no correction (11 of 24 cells, mostly GB) and increases the mean absolute calibration intercept in every cell. Conclusions are limited to one dataset and 10 seeds.
